\documentclass[11pt]{article}
\usepackage{fontspec}
\usepackage[french,english]{babel}
\usepackage{amsmath}
\usepackage{amssymb}
\usepackage{amsthm}
\usepackage{mathrsfs}
\usepackage{tikz-cd}
\usepackage{enumitem}
\usetikzlibrary{babel}

\newcommand{\printedpage}[1]{\marginpar{\footnotesize p.~#1}}
\newcommand{\msp}[1]{\marginpar{\footnotesize [#1]}}
% dotted arrow as printed in "c : x_0 ....> x_1"
\newcommand{\dotto}{\mathrel{\tikz[baseline=-0.5ex]\draw[dotted,-latex] (0,0)--(1.6em,0);}}

\begin{document}

\printedpage{10}
% head: A. Grothendieck : Dérivateurs, Ch. XII, Édition des Universités de Montpellier II et Paris VII
% foot: Version 20 juin 2003
\textsc{Proof}. Let
\[ F \hookrightarrow U \]
be the inclusion into $U = \partial X$ of the corner-category \texttt{[}$V =$\texttt{]} $\vcenter{\hbox{\begin{tikzcd}[row sep=small, column sep=small]
0 \arrow[r] \arrow[d] & a \\
b &
\end{tikzcd}}}$ (cf. figure), isomorphic to $\Psi$. One sees at once that this is a closed Dwyer immersion (the smallest open of $U$ containing $\Psi$ is just the category $\tilde{\Psi} = U/t_1$ from earlier), hence a $W$-cofibration. If $U \longrightarrow X$ were a $W$-cofibration, the same would hold for the composite $V \longrightarrow X$. We are going to see that this is not the case.

Let $\overline{V} = \vcenter{\hbox{\begin{tikzcd}[row sep=small, column sep=small]
0 \arrow[r] \arrow[d] & a \arrow[d] \\
b \arrow[r] & t_1
\end{tikzcd}}}$ be the closure of $V$ in $X$,
\msp{page 14}and consider the inclusion
\[ \overline{V} \longrightarrow X, \]
I say that \texttt{[it]} is a closed Dwyer immersion, i.e. that for every $s \in W\backslash\overline{V}$, where $W$ is the smallest open of $X$ containing $\overline{V}$, i.e.
\[
W = \quad
\begin{tikzcd}[row sep=small, column sep=small]
0 \arrow[r] \arrow[d] & a \arrow[d] & a' \arrow[l] \arrow[d] \\
b \arrow[r] & {} & a'' \arrow[l] \\
b' \arrow[u] \arrow[r] & b'' \arrow[u] & c, \arrow[l] \arrow[u]
\end{tikzcd}
\]
the category $s\backslash\overline{V}$ has an initial object. (The categories $s\backslash\overline{V}$ that one obtains, from the five possible choices for $s$, are $\boldsymbol{\Delta}^1$ (if $s = a'$ or $b'$) or $e$ (for $s = a''$, $b''$ or $c$)). Thus $\overline{V} \longrightarrow X$ is a $W$-cofibration.

This being established, one knows that there exists, for [?] the inclusion (open immersion)
\[ V \hookrightarrow \overline{V} \qquad \text{i.e. } \Psi \hookrightarrow \boldsymbol{\Delta}^1 \times \boldsymbol{\Delta}^1, \]
a functor $V \longrightarrow V'$, giving rise to a cocartesian square
\[
\begin{tikzcd}
V \arrow[r, hook, "i"] \arrow[d, "p"'] & \overline{V} \arrow[d, "\overline{p}"] \\
V' \arrow[r, hook, "i'"] & \overline{V'}
\end{tikzcd}
\tag*{$(*)$}
\]
with
\[ i, p \in W \ \text{(for $V$, $V'$, $\overline{V}$ $W$-aspheric), but } i', \overline{p} \notin W \ \text{(i.e. $\overline{V'}$ \textit{not} $W$-aspheric).} \]

I recall that one first takes
\printedpage{11}
% head: A. Grothendieck : Dérivateurs, Ch. XII, Édition des Universités de Montpellier II et Paris VII
% foot: Version 20 juin 2003
\[
V'_0 = \left(
\begin{tikzcd}[row sep=small]
 & & a \\
0' \arrow[r, shift left, "\lambda"] \arrow[r, shift right, "\mu"'] & 0 \arrow[ur, "\alpha"] \arrow[dr, "\beta"'] & \\
 & & b
\end{tikzcd}
\right)
\qquad \text{with} \quad
\left\{
\begin{aligned}
\alpha\lambda &= \alpha\mu \\
\beta\lambda &= \beta\mu
\end{aligned}
\right.
\]
and the ``identical'' closed immersion of $V$ into $V'$ \texttt{[rather} $V'_0$\texttt{]}
\[ p_0 : V \longrightarrow V'_0 \]
\msp{page 15}(which is such that
\[ p_{0!}(e_V) = e_{V'_0} \ \ ), \]
then for $V'$ one takes the category obtained by adding a final object to $V'_0$
\[
V' = \left(
\begin{tikzcd}[row sep=small, column sep=small]
 & & a \arrow[dr, "\rho"] & \\
0' \arrow[r, shift left, "\lambda"] \arrow[r, shift right, "\mu"'] & 0 \arrow[ur, "\alpha"] \arrow[dr, "\beta"'] & & e \\
 & & b \arrow[ur, "\sigma"'] &
\end{tikzcd}
\right)
\qquad \text{with} \quad
\left\{
\begin{aligned}
\alpha\lambda &= \alpha\mu \\
\beta\lambda &= \beta\mu \\
\rho\alpha &= \sigma\beta
\end{aligned}
\right.
\]
and for $p$ the composite functor
\[
\begin{tikzcd}
V \arrow[r, "p_0"] \arrow[rr, bend right, "p"'] & V'_0 \arrow[r, "p_1"] & V',
\end{tikzcd}
\]
where $p_1$ is the inclusion (open). One then finds the square $(*)$, with $\overline{V'}$ not $W$-aspheric, or what amounts to the same thing, $i' \notin W$. For since $\overline{V'}\backslash V' \simeq \overline{V}\backslash V$ is the one-point category, $i' \in W$ is equivalent to $V'/t$ $W$-aspheric (\footnote[8]{assuming (W 7 bis).}) (cf. cor. 4, p. 149 in XII). Now by construction $V'/t \simeq V'_0$, which is not $W$-aspheric (in the case, say, where $W \subseteq W_1$, for $\pi_1(V'_0) \simeq \mathbf{Z}$).

\msp{page 16}Let us then consider the diagram of cocartesian squares
\[
\begin{tikzcd}
V \arrow[r, hook, "i"] \arrow[d, "p"'] & \overline{V} \arrow[r, hook] \arrow[d, "\overline{p}"] & X \arrow[d, "q"] \\
V' \arrow[r, hook, "i'"] & \overline{V'} \arrow[r, hook] & X'
\end{tikzcd}
\qquad .
\]

\printedpage{12}
% head: A. Grothendieck : Dérivateurs, Ch. XII, Édition des Universités de Montpellier II et Paris VII
% foot: Version 20 juin 2003
Since $\overline{V} \hookrightarrow X$ is \texttt{[a]} closed Dwyer immersion, the square on the right is $W$-cocartesian, hence $(\overline{p} \in W) \iff (q \in W)$, hence (since $\overline{p} \notin W$) we have
\[ p \in W, \qquad \text{but} \qquad q \notin W, \]
which indeed shows that $V \longrightarrow X$ is not a $W$-cofibration. Factoring $V \longrightarrow X$ through $U$ (instead of $\overline{V}$), one finds
\[
\begin{tikzcd}
V \arrow[r] \arrow[d, "p\in W"'] & U \arrow[r] \arrow[d, "q_U\in W"] & X \arrow[d, "q\notin W"] \\
V' \arrow[r] & U' \arrow[r] & X'
\end{tikzcd}
\qquad ,
\]
where $q_U \in W$ since $p \in W$ and $V \longrightarrow U$ is \texttt{[a]} Dwyer immersion. This makes explicit that $U \longrightarrow X$ is not a $W$-cofibration, since
\[ q_U \in W, \qquad \text{but} \qquad q \notin W, \]
\hfill q.e.d.

\msp{page 17}\textbf{Gluing in} Cat \textbf{of an open and a complementary closed.}

\texttt{[Suppose given]}
\[ X_0 \underbrace{\hookrightarrow}_{\text{open}} X \ \underbrace{\hookleftarrow}_{\substack{\text{closed}\\ \text{comple-}\\ \text{mentary}}} X_1 \]
in Cat, \texttt{[and let]}
\[ \tilde{X}_{01} \subseteq \mathrm{Fl}(X) \]
be \texttt{[the]} full subcategory of arrows with source in $X_0$, target in $X_1$.

\textbf{Question.} \textit{This diagram}
\[
\begin{tikzcd}[column sep=small]
 & & \tilde{X}_{01} \arrow[dll, "p"' near start, "\text{fibrant, fibers the $x_0\backslash X_1$}"' near end] \arrow[drr, "q" near start, "\text{cofibrant, fibers the $X_0/x_1$}" near end] & & \\
X_0 \arrow[ddrr, "i"'] & & & & X_1 \arrow[ddll, "j"] \\
 & {} \arrow[rr, bend left] & & {} & \\
 & & X, & &
\end{tikzcd}
\tag*{$(*)$}
\]
\printedpage{13}
% head: A. Grothendieck : Dérivateurs, Ch. XII, Édition des Universités de Montpellier II et Paris VII
% foot: Version 20 juin 2003
\textit{is it $W$-cocartesian?}

OK if $i$ $W$-aspheric or j $W$-coaspheric, i.e. $X \longrightarrow \boldsymbol{\Delta}^1$ smooth or proper -- this is a case where one has either $X_0 \xrightarrow{W} X$, or $X_1 \xrightarrow{W} X$. But even assuming $i$ or $j$ in $W$, and moreover $X$ aspheric (hence $X_0$ and $X$ aspheric, or $X_1$ and $X$ aspheric), and $W = \mathbf{W}_\infty$, I have not been able to prove that the diagram is cocartesian, i.e. $q \in W$ resp. $p \in W$ (\footnote[9]{I assume that $W$ satisfies Loc (6).}). (Nor if $X_0 \simeq e$, $X$, $X_1$ $W$-aspheric $\Longrightarrow$ $e\backslash X$ $W$-aspheric, i.e. $X_1 \longrightarrow X$ $W$-coaspheric $\ldots$\,) But I have at least proved, in all cases, that % [sic] j (upright) -> $j$, letter denoting a functor printed in roman type
\[
\begin{tikzcd}[row sep=small]
 & X_0 & & \\
\mathcal{X} \overset{\text{def}}{=} \displaystyle\int \tilde{X}_{01} \arrow[ur] \arrow[dr] & & {} \arrow[r, "\varphi"] & X \\
 & X_1 & &
\end{tikzcd}
\tag*{$(**)$} % verified: CMSY10 glyph, \mathcal{X}; corrections.csv, PDF leaf 14
\]
is in $\mathbf{W}_1$.

Something \textit{false} (unless $i$ or $j$ in $W = \mathbf{W}_\infty$) if one replaces $\tilde{X}_{01}$ by $\tilde{X}'_{01}$, formed of the triples $(x_0, x_1, c)$, with $x_0 \in \mathrm{Ob}\,X_0$, $x_1 \in \mathrm{Ob}\,X_1$, $c : x_0 \dotto x_1$ a \textit{path} from $x_0$ to $x_1$ in $X$! We have
\msp{page 18}therefore a ``topological'' situation entirely peculiar to the category of models Cat, which does not seem to have an analog in categories such as Esp or $\boldsymbol{\Delta}^\wedge$, where however the notions of open and complementary closed (with a grain of salt in $\boldsymbol{\Delta}^\wedge$) have a meaning.

The quasi-trivial categorical argument which establishes that $(**)$ is in $\mathbf{W}_1$, i.e. induces an equivalence of categories between $\mathrm{Rev}(X)$ and
\[ \mathrm{Rev}(\mathcal{X}) \ \simeq \ \mathrm{Rev}(X_0) \overset{2}{\times}_{\mathrm{Rev}(\tilde{X}_{01})} \mathrm{Rev}(X_1), \]
seems to me to have to extend to higher dimensions, by dint of $n$-gerbes on $X$ -- but the foundations for the theory of $n$-gerbes are lacking. (And I have not had the courage to spell out the argument even for $n = 2$ (where however \textsc{Giraud} provides what is needed) (\footnote[10]{But I have nevertheless made a sketch of a verification for $n = 2$, which convinces me that we do indeed have $\varphi \in \mathbf{W}_2$. And I do not have hardly any doubt that we even have $\varphi \in \mathbf{W}_\infty$.}).) So I conjecture that $\varphi$ is at the very least in $\mathbf{W}_\infty$, i.e. that the diagram $(*)$ is $W$-cocartesian if $W \supseteq \mathbf{W}_\infty$ (and satisfies Loc (6 a$'$)). % [sic] je n'ai pas guère (footnote 10) -> je n'ai guère, double negation pas + guère

\end{document}
